\documentclass[a4paper, serif, 10pt]{moderncv}

%% moderncv
% color and style
\moderncvcolor{black}
\moderncvstyle{banking}

%% page layout/margins
\usepackage[top=2.5cm, bottom=2.5cm, left=1.5cm, right=1.5cm]{geometry}

\nopagenumbers{}

%% Babel config for Brazilian Portuguese language
% ref:
% - https://latex3.github.io/babel/guides/locale-portuguese.html
\usepackage[brazilian]{babel}

%% fontspec
\usepackage{fontspec}

\IfFontExistsTF{Linux Libertine}{
  \setmainfont[Ligatures={TeX, Common}, Numbers=OldStyle]{Linux Libertine}
}{}
\IfFontExistsTF{Linux Libertine O}{
  \setmainfont[Ligatures={TeX, Common}, Numbers=OldStyle]{Linux Libertine O}
}{}

%% CTeX
% CJK support, used to show CJK characters in the resume
%
% - heading=false: disable ctex heading
% - punct=kaiming: use kaiming punctuations styles for CJK
% - scheme=plain: use plain scheme, do not override \normalsize font size
% - space=auto: space settings for CJK characters
%
% CJK fonts are resolved per locale language, e.g. Noto Serif/Sans CJK JP
% for Japanese resumes, with fallback to the next region if not installed
%
% ref:
% - http://ctan.mirrorcatalogs.com/language/chinese/ctex/ctex.pdf
\usepackage[UTF8, heading=false, punct=kaiming, scheme=plain, space=auto]{ctex}

\IfFontExistsTF{Noto Serif CJK SC}{
  \setCJKmainfont{Noto Serif CJK SC}
}{}
\IfFontExistsTF{Noto Sans CJK SC}{
  \setCJKsansfont{Noto Sans CJK SC}
}{}

%% line spacing
\usepackage{setspace}
\setstretch{1.125}

%% URL styling - use normal text instead of monospace
\urlstyle{same}

%% Patch \httplink and \httpslink to handle full URLs
% The original moderncv macros always prepend http:// or https:// to URLs.
% This causes issues when the URL already has a protocol (e.g., https://example.com)
% resulting in malformed URLs like https://https://example.com.
% This patch checks if the URL already contains "://" and uses it directly if so.
% Note: We use \str_set:Nx (x-expansion) to expand macros like \@homepage first.
\ExplSyntaxOn
\str_new:N \g_yamlresume_url_str

\RenewDocumentCommand{\httpslink}{O{}m}{
  \str_set:Nx \g_yamlresume_url_str {#2}
  \str_if_in:NnTF \g_yamlresume_url_str {://}
    { \url{#2} }
    { \url{https://#2} }
}

\RenewDocumentCommand{\httplink}{O{}m}{
  \str_set:Nx \g_yamlresume_url_str {#2}
  \str_if_in:NnTF \g_yamlresume_url_str {://}
    { \url{#2} }
    { \url{http://#2} }
}
\ExplSyntaxOff

\name{Mariana Costa Almeida}{}
\title{Analista de Negócios Sênior | Transformação Digital e Dados}
\phone[mobile]{+55 11 98765-4321}
\email{mariana.almeida@exemplo.com.br}
\homepage{https://www.marianaalmeida.com.br}

\address{Av. Paulista, 1578, Bela Vista, São Paulo, SP, Brasil, 01310-200}{}{}

\extrainfo{{\small \faLinkedin}\ \href{https://www.linkedin.com/in/mariana-costa-almeida}{@mariana-costa-almeida} {} {} {} • {} {} {} 
{\small \faMedium}\ \href{https://medium.com/@marianacalmeida}{@marianacalmeida}}

\begin{document}

\maketitle

\section{Informações Básicas}

\cvline{}{Analista de negócios sênior com mais de 8 anos de experiência em projetos de transformação digital, melhoria de processos e produtos digitais nos setores financeiro e de serviços.
%
\begin{itemize}
\item Especialista em levantamento e priorização de requisitos, modelagem de processos (BPMN) e tradução de necessidades de negócio em soluções técnica e comercialmente viáveis.
\item Experiência com análise de dados, definição de indicadores-chave e painéis executivos em SQL e Power BI para apoiar decisões e acompanhar resultados.
\item Forte atuação com partes interessadas multidisciplinares, facilitação de oficinas, gestão de backlog e metodologias ágeis (Scrum e Kanban).
\end{itemize}}
\section{Formação}

\cventry{fev. de 2012–dez. de 2016}
        {Bacharelado, Administração de Empresas, Nota: 8.7/10}
        {Universidade de São Paulo}
        {\url{https://www.usp.br}}
        {}
        {\begin{itemize}
\item Formação sólida em gestão, processos e tecnologia, com ênfase em análise de negócios e melhoria contínua.
\item Desenvolvimento de projetos acadêmicos de mapeamento de processos e avaliação de indicadores de desempenho.
\end{itemize}
\textbf{Cursos}: Economia de Empresas, Estatística Aplicada, Gestão de Processos, Sistemas de Informação, Finanças Corporativas, Pesquisa Operacional}

\cventry{jan. de 2018–dez. de 2019}
        {Mestrado, Gestão de Projetos e Inteligência de Negócios, Nota: 9.1/10}
        {Fundação Getulio Vargas}
        {\url{https://www.fgv.br}}
        {}
        {\begin{itemize}
\item Aprofundamento em gestão de projetos, governança de dados e inteligência de negócios aplicados à tomada de decisão.
\end{itemize}
\textbf{Cursos}: Gestão de Projetos, Governança de Dados, Inteligência de Negócios, Análise de Decisão, Metodologias Ágeis}
\section{Experiência Profissional}

\cventry{mar. de 2021–Atual}
        {Analista de Negócios Sênior}
        {Nexo Soluções Digitais}
        {\url{https://www.nexosolucoes.com.br}}
        {}
        {\begin{itemize}
\item Liderança de iniciativas de transformação digital para áreas de crédito, cobrança e atendimento, com foco em eficiência operacional e experiência do cliente.
\item Condução de oficinas de descoberta, entrevistas com usuários e análise de dados para priorizar requisitos e construir roteiros alinhados aos objetivos de negócio.
\item Especificação de histórias de usuário, critérios de aceite e fluxos BPMN, apoiando times de desenvolvimento e QA durante todo o ciclo de entrega.
\item Acompanhamento de indicadores-chave e elaboração de painéis em Power BI para reporte executivo e identificação de oportunidades de melhoria.
\end{itemize}
\textbf{Palavras-chave}: Levantamento de requisitos, Modelagem de processos, BPMN, Power BI, SQL, Gestão de partes interessadas}

\cventry{jan. de 2018–fev. de 2021}
        {Analista de Negócios}
        {Banco Horizonte}
        {\url{https://www.bancohorizonte.com.br}}
        {}
        {\begin{itemize}
\item Atuação na análise e melhoria de processos de entrada de clientes, reduzindo o tempo médio de abertura de conta e retrabalho operacional.
\item Levantamento de requisitos funcionais e não funcionais junto às áreas de negócio, conformidade e tecnologia.
\item Apoio na implantação de práticas ágeis, gestão de backlog e documentação de regras de negócio em Confluence.
\item Realização de testes de aceitação e acompanhamento de indicadores de qualidade das entregas.
\end{itemize}
\textbf{Palavras-chave}: Análise de dados, Metodologias ágeis, Melhoria contínua, Documentação funcional, Testes de aceitação}

\cventry{jan. de 2017–dez. de 2017}
        {Analista de Processos Júnior}
        {Consultoria Ponte Sul}
        {\url{https://www.consultoriapontesul.com.br}}
        {}
        {\begin{itemize}
\item Mapeamento de processos AS-IS e proposição de fluxos TO-BE para rotinas administrativas e financeiras.
\item Elaboração de documentação funcional, casos de uso e relatórios de acompanhamento para clientes do setor de serviços.
\item Suporte à implantação de melhorias e treinamento de usuários em novas rotinas e ferramentas.
\end{itemize}
\textbf{Palavras-chave}: AS-IS/TO-BE, Casos de uso, Processos, Treinamento}
\section{Idiomas}

\cvline{Português}{Nativo ou Bilíngue \hfill \textbf{Palavras-chave}: Nativo}
\cvline{Inglês}{Proficiência Profissional Fluente \hfill \textbf{Palavras-chave}: TOEIC 920}
\cvline{Espanhol}{Proficiência Profissional Limitada \hfill \textbf{Palavras-chave}: Leitura técnica}
\section{Competências}

\cvline{Análise de Negócios}{Especialista \hfill \textbf{Palavras-chave}: BPMN, Histórias de usuário, Casos de uso, AS-IS/TO-BE, Gestão de partes interessadas}
\cvline{Dados e BI}{Avançado \hfill \textbf{Palavras-chave}: SQL, Power BI, Excel avançado, Tableau, Indicadores}
\cvline{Metodologias Ágeis}{Avançado \hfill \textbf{Palavras-chave}: Scrum, Kanban, Backlog, Jira}
\cvline{Ferramentas}{Avançado \hfill \textbf{Palavras-chave}: Jira, Confluence, Visio, Figma, Google Workspace}
\section{Prêmios}

\cventry{nov. de 2022}
        {Nexo Soluções Digitais}
        {Prêmio de Excelência em Projetos}
        {}
        {}
        {Reconhecimento pelo desempenho na condução de projetos estratégicos de transformação digital, com destaque para a redução de custos operacionais e a melhoria da satisfação dos clientes.}

\cventry{dez. de 2016}
        {Universidade de São Paulo}
        {Destaque Acadêmico}
        {}
        {}
        {Reconhecimento acadêmico por desempenho de destaque no curso de Administração de Empresas.}
\section{Certificados}

\cventry{ago. de 2023}
        {IIBA}
        {Certified Business Analysis Professional (CBAP)}
        {\url{https://www.iiba.org/certification/cbap/}}
        {}
        {}

\cventry{mai. de 2022}
        {Scrum.org}
        {Professional Scrum Product Owner I}
        {\url{https://www.scrum.org/professional-scrum-product-owner-certifications}}
        {}
        {}

\cventry{set. de 2021}
        {Microsoft}
        {Microsoft Certified: Power BI Data Analyst Associate}
        {\url{https://learn.microsoft.com/en-us/certifications/power-bi-data-analyst-associate/}}
        {}
        {}
\section{Publicações}

\cventry{jun. de 2023}
        {Gestão de requisitos em projetos de transformação digital}
        {Revista Brasileira de Gestão e Inovação}
        {\url{https://www.revistabgi.com.br/artigos/gestao-de-requisitos-em-projetos-de-transformacao-digital}}
        {}
        {\begin{itemize}
\item Artigo sobre práticas de levantamento, priorização e rastreabilidade de requisitos em iniciativas de transformação digital.
\item Discussão de indicadores de sucesso e lições aprendidas em projetos dos setores financeiro e de serviços.
\end{itemize}}
\section{Projetos}

\cventry{abr. de 2022–out. de 2022}
        {Redesenho da jornada de entrada de clientes, com pesquisa qualitativa, mapeamento de pontos de dor e priorização de melhorias.}
        {Mapeamento da Jornada do Cliente}
        {}
        {}
        {\begin{itemize}
\item Condução de entrevistas e oficinas com áreas de negócio, atendimento e tecnologia.
\item Mapeamento da jornada atual e proposição de melhorias com foco em autoatendimento e redução de etapas manuais.
\item Definição de indicadores de sucesso e acompanhamento dos resultados após a implantação.
\end{itemize}
\textbf{Palavras-chave}: Jornada do cliente, Desenho de serviços, BPMN}

\cventry{fev. de 2023–jul. de 2023}
        {Painel gerencial para acompanhamento de indicadores operacionais e financeiros, utilizado pela diretoria na tomada de decisão.}
        {Painel Executivo de Indicadores}
        {}
        {}
        {\begin{itemize}
\item Levantamento de requisitos e definição de indicadores-chave junto às áreas de negócio.
\item Modelagem de dados em SQL e construção de visualizações no Power BI.
\item Documentação de regras de negócio e treinamento de usuários-chave.
\end{itemize}
\textbf{Palavras-chave}: Power BI, Indicadores-chave, SQL, Governança de dados}
\section{Interesses}

\cvline{Educação e Mentoria}{Mentoria, Palestras, Cursos}
\cvline{Bem-estar}{Corrida, Yoga, Culinária}
\section{Voluntariado}

\cventry{mar. de 2022–Atual}
        {Mentora de Jovens Analistas}
        {Instituto Ponte para o Futuro}
        {\url{https://www.ponteparaofuturo.org.br}}
        {}
        {\begin{itemize}
\item Mentoria de jovens em início de carreira em análise de negócios, com orientação sobre requisitos, processos e comunicação com partes interessadas.
\item Participação em encontros de carreira e revisão de currículos para candidatos de baixa renda.
\end{itemize}}
    
\end{document}